\begin{abstract}
Neural receivers (NRXs) built from convolutional and graph neural network (CGNN) layers obtain their accuracy from iterative message passing, but in the real-time implementation of a state-of-the-art NRX for the 5G NR physical uplink shared channel (PUSCH), a 1\,ms latency budget admits only two of its eight iterations. We propose a stateful NRX that keeps a pilot memory, i.e., the least-squares channel estimates at the pilot positions of earlier slots of each user. A channel filter combines them with the current pilots; its weights are not trained but fitted per user by least squares on the memory itself, which requires neither channel statistics nor labels and adds no weights to the NRX. In $2\times 4$ multi-user MIMO simulations with QPSK, 16-QAM, and 64-QAM, the proposed NRX with two iterations requires 0.5 to 0.9\,dB less SNR than a stateless two-iteration NRX retrained on the same data and outperforms the released three-iteration NRX in all cases, with and without power control; with QPSK, it also outperforms the released eight-iteration NRX in all cases. It keeps most of its memory gain with phase jumps, frequency offsets, and TDD uplink slots, and it meets the 1\,ms real-time budget, which the three-iteration NRX exceeds.
\end{abstract}
